\chapter{Personal Identity and the Ship of Theseus}
\label{ch:ship_of_theseus}

\papernote{Draws from: \emph{Unified Entropy from Oscillation, Category, and Partition} --- \texttt{blindhorse/docs/categorical-partitioning/} (Ship of Theseus section); and Chapter~\ref{ch:identity} (non-grounded naming circuits)}

\begin{chapterabstract}
The Ship of Theseus---if every plank of a ship is replaced over time, is it the same ship?---is resolved within the categorical framework as a partition continuity problem. Identity is not spatial continuity but categorical trajectory continuity: the same sequence of completed categorical states, connected by the non-grounded naming circuit. A system retains identity so long as its naming circuit maintains continuity of the completion sequence, regardless of material replacement. This dissolves not only the Ship of Theseus but also the sorites (heap) paradox and Zeno's paradoxes, all of which are shown to be instances of partition boundary ambiguity at finite resolution.
\end{chapterabstract}

\input{../blindhorse/docs/categorical-partitioning/sections/ship-of-theseus.tex}
\input{../blindhorse/docs/categorical-partitioning/sections/heap-paradox.tex}
\input{../blindhorse/docs/categorical-partitioning/sections/zenos-paradox.tex}
